\section{Audit and organize the existing databases}
\label{sec:data}
\subsection{Your resources and the first decisions}
The user confirmed an existing detailed COMSOL database and access to the real measurements from Bathla et al.~\cite{bathla}. The experimental files are on external storage and will be obtained later; their contents have not yet been inspected. The first-paper output is confirmed as near-plane magnetic fields, with currents used as an auxiliary model. Begin the simulation study immediately; make the real-data protocols conditional on the audit rather than postponing the entire project.

\subsection{COMSOL inventory: the first two working days}
Make a read-only manifest before conversion. Give each source layout a stable \texttt{geometry\_id}, and each simulation configuration a \texttt{scene\_id}. Distinguish a genuinely new conductor layout from a different height, current, mesh, or crop of the same layout. Count unique geometries, not exported images.

\begin{longtable}{P{37mm}P{114mm}}
\toprule Required field & Purpose and verification\\\midrule\endhead
Geometry/source identity & CAD/model hash; topology; turns; width; spacing; source layer; contacts. Identify duplicates and translated copies.\\
Coordinates and pitch & Physical $x,y,z$ coordinates; ascending order; exact $\Delta x,\Delta y$; FOV extent. Verify orientation and registration.\\
Height convention & Origin of $z$; conductor top, centre, or substrate reference; known source thickness; absolute height versus gap.\\
Drive & Applied current/voltage; measured or prescribed value; sign; contact boundary conditions.\\
Physics and quantity & $\mathbf B$ versus $\mathbf H$; units; material permeability; current density availability; source-free region validity.\\
Numerics & COMSOL version, mesh size, solver tolerance, interpolation/export procedure, simulation boundaries.\\
Files and provenance & Paths, SHA-256, grid shape, data type, component labels, missing values, source-generation seed.\\
Target availability & Field at $z_{\rm ref}$; existence of complete same-geometry height sweeps; whether a trusted target needs new simulation.\\\bottomrule
\end{longtable}

Create a report showing height counts per geometry, pitch/FOV distributions, missing components, field ranges, and coordinate agreement. Plot a few profiles at every height. Repeated field interpolation artifacts are model discrepancy, not automatically independent measurement noise.

\subsection{Reuse complete height sweeps without learning test-pair kernels}
Where a geometry has fields at $z_{\rm ref}$ and several higher planes, use the lower field as target and each higher plane as an observation. The analytical operator is evaluated from physical spacing and height. A held-out sharp field must not be used to fit that observation's empirical kernel.

Compare direct COMSOL high-plane fields with analytical propagation of the low-plane field. This residual includes export interpolation, discretization, finite-window error, registration, and any modelling mismatch. Investigate its structure before ascribing it to mesh noise. A bounded empirical correction may be a later experiment, but its parameters must be fitted on training/calibration data only, and tested on unused geometries or acquisition sessions.

\subsection{Dataset tiers}
\begin{tabularx}{\linewidth}{P{27mm}YY}
\toprule Tier & Data & Scientific role\\\midrule
T0: mechanics & Small analytical cases, Fourier modes, closed loops & Operator, adjoint, gradients, and ambiguity checks\\
T1: pilot & Stratified subset of the existing COMSOL database & Baseline comparison and failure discovery\\
T2: challenge & Disjoint COMSOL layouts/families and selected numerical settings & Generalization; simulator/model mismatch\\
T3: real archive & Audited Bathla maps, labels, and repetitions & Real noise, dictionary comparison, source activations\\
T4: new acquisition & Calibrated height sweeps and independent layouts & Direct validation of target-plane reconstruction and distance\\\bottomrule
\end{tabularx}
Use new analytical scenes to answer controlled questions, not as a replacement for the existing database. Add new COMSOL simulations only to fill an identified gap: geometry families, near-plane references, mesh-convergence checks, or height ranges.

\subsection{Splitting protocol and leakage controls}
First assign source identities to train, validation, calibration, and test groups; only then derive heights, current rescalings, crops, and noise variants. The calibration group is needed only if uncertainty calibration is included and must be distinct from the model-selection validation set. All descendants of a geometry remain in the same group. Near-duplicates and variants from the same template generator may require a higher-level family grouping.

Maintain two different evaluations: \textbf{new-instance interpolation} within a geometry family, and \textbf{unseen-family transfer}. A random split of configurations is not evidence of unseen geometry. Separate height interpolation from extrapolation. Record exact withheld ranges and test them once after selecting the model.

\begin{figure}[htbp]\centering
\begin{tikzpicture}[node distance=7mm]
\node[wideblock] (ids) {Independent source/CAD identities\\and acquisition sessions};
\node[block,below left=13mm and 22mm of ids] (train) {Training sources};
\node[block,below=13mm of ids] (val) {Validation and\\calibration sources};
\node[block,below right=13mm and 22mm of ids] (test) {Locked test sources};
\node[block,below=of train] (td) {All heights, crops,\\noise descendants};
\node[block,below=of val] (vd) {All heights, crops,\\noise descendants};
\node[block,below=of test] (ed) {All heights, crops,\\noise descendants};
\draw[arrow] (ids)--(train);\draw[arrow] (ids)--(val);\draw[arrow] (ids)--(test);
\draw[arrow] (train)--(td);\draw[arrow] (val)--(vd);\draw[arrow] (test)--(ed);
\end{tikzpicture}
\caption{Split physical identities before creating training examples. Correlated crops and height variants cannot be moved across source-level partitions.}
\end{figure}

For a pilot, choose approximately 100 independent layouts \emph{if available}, with a 60/20/20 source-level train/validation/test split. This is a starting engineering configuration, not a sample-size justification. If calibration is needed, reserve additional sources or repartition without examining the locked test outcomes. Expand according to learning curves and precision of source-level confidence intervals.

\subsection{Sampling, cropping, and normalization}
Propagate on a larger domain and then apply pixel integration and FOV restriction. For coarser pitch, antialias and integrate the field; do not simply resize an image and declare a new physical spacing. Upsampling labels does not create new measured information. If the target uses a different grid, make that output grid explicit and interpolate only through a documented observation model.

Use a fixed scale estimated from training data or an observation-only scale $q(y)$, and retain $q$ as physical metadata when amplitude carries information. Reconstruct in physical units for metrics. If minimizing the existing half-squared loss plus an L1 penalty under $x=q u$, the normalized coefficient is $\lambda/q$. Do not normalize using held-out sharp targets. At zero or near-zero signal, use a predetermined scale floor and flag height as uninformative.

Rotating a field image requires rotating its vector components as well as its spatial coordinates. Under a planar rotation $R$, rotate $(B_x,B_y)$ consistently and transform the NV projection direction for scalar observations. Scaling geometry changes physics and cannot be implemented by spatial resizing alone. Reverse all signed components for current reversal under the linear magnetostatic model.

\section{How to use Anuj Bathla's experimental dataset}
\label{sec:bathla}
\subsection{What is established and what is still unknown}
The published study describes measured individual-coil basis fields and mixed activations on a multilayer micro-coil platform, plus time-dependent measurements. Its reconstruction uses sparse coefficients in a measured dictionary. Your group's original files may contain more information than the paper, but this guide cannot assume that fixed-reference sharp fields, absolute calibrated standoffs, or all vector components are already present. Obtain those facts from acquisition logs.

\subsection{Data-receipt checklist}
\begin{enumerate}
\item Copy data into a read-only archive; retain original folder structure, checksums, and the acquisition notebook. Work from versioned derived copies.
\item Identify field maps, raw ODMR/lock-in sequences, current-on/off references, individual-coil maps, mixed activations, timestamps, and event waveforms.
\item Obtain the coil-to-layer lookup, microscope magnification, pixel pitch, NV axis, coordinate conventions, drive values, and registration transforms.
\item Label every distance as independently measured, nominal, model-fitted, relative stage displacement, or unknown. Store uncertainty and calibration method.
\item Ask whether basis maps and mixtures were measured in the same session and at the same drive. Check drift, gain, and background differences.
\item Verify data-release permission and what can be made public. This affects the reproducibility package; it does not alter the scientific split.
\end{enumerate}

\subsection{Six concrete uses, in increasing order of difficulty}
\textbf{U1: benchmark the published task.} Reproduce field fitting and activation recovery with the measured dictionary before adding a neural network. The unknown is a coefficient vector, not field pixels. For dictionary $D=[d_1,\ldots,d_K]$ and observation $y$, solve
\begin{equation}
 \hat c=\arg\min_c \tfrac12\norm{W(Dc-y)}^2+\lambda\norm{c}_1.
\end{equation}
Use NNLS only if nonnegative currents are physically justified. Compare coefficient error, activation precision/recall, and field residuals. This is a calibration and reproducibility exercise, not a new research claim.

\textbf{U2: learn experimental noise.} Repetitions and current-off acquisitions can estimate spatial variance, correlations, drift, and fitting uncertainty. Subtracting a spatially varying physical field is not equivalent to measuring noise. Split sessions so repeated backgrounds do not leak. Where raw ODMR fits provide covariance, propagate it or use it to weight the observation loss.

\textbf{U3: calibrate simulator-to-measurement discrepancy.} Match known CAD and drive settings to measured individual-coil fields. Fit limited nuisance parameters such as global registration, axis angle, gain, and effective distance on a calibration subset. Do not allow a free learned kernel to absorb every error. Test fitted corrections on unused coils or sessions.

\textbf{U4: experimentally adapt a pretrained field prior.} Start from COMSOL supervision. Use repeated real observations, known source activations, and correctly modelled measurement consistency to adapt on an explicit adaptation set. Validate on separate acquisitions. A dictionary-generated $D c$ is a \emph{model-based surrogate at the sensor plane}, not an independent near-plane truth. Shared noisy basis maps can correlate surrogate target and input noise.

\textbf{U5: construct a real dictionary LISTA baseline.} Train an unrolled sparse-code network for $c$ on activation labels and compare with LASSO at matched inference latency. This tests learned acceleration and provides a source-aware auxiliary branch. It does not directly establish standoff super-resolution. Hold out activation patterns for known-dictionary performance; hold out source identities or a new array for geometry transfer.

\textbf{U6: validate near-plane reconstruction.} If the archive contains registered measurements at several calibrated positions, form higher-to-lower-plane test pairs with the same optics. Otherwise acquire new height sweeps. Do not invent sharp references by applying your own inverse to the archive and then evaluate against them.

\begin{figure}[htbp]\centering
\begin{tikzpicture}[node distance=7mm]
\node[wideblock] (archive) {Bathla experimental archive\\audit labels and calibration};
\node[wideblock,below left=11mm and 8mm of archive] (known) {Basis maps + mixtures\\dictionary/activation benchmark};
\node[wideblock,below right=11mm and 8mm of archive] (repeat) {Raw/repeated acquisitions\\noise and domain adaptation};
\node[wideblock,below=of known] (labels) {Coefficients and sensor-plane fields\\are useful supervision};
\node[wideblock,below=of repeat] (sweep) {Calibrated height pairs\\validate near-plane prediction};
\draw[arrow] (archive)--(known);\draw[arrow] (archive)--(repeat);
\draw[arrow] (known)--(labels);\draw[arrow] (repeat)--(sweep);
\end{tikzpicture}
\caption{Different experimental data support different claims. Activate a branch only after the necessary metadata and measurements are verified.}
\end{figure}

\subsection{Splits that are honest for a single fabricated array}
Holding out activation combinations tests combinatorial generalization for \emph{known source templates}. All combinations on one array are not independent new geometries. A shared dictionary may be legitimate calibration information in that task, but it must be declared and kept separate from unseen-geometry claims. Use held-out acquisition sessions to test drift. For dynamics, split contiguous events/runs and leave guard intervals; neighbouring frames are strongly correlated.

If a known-coil dictionary is available at test time, report that regime separately from ``no per-coil calibration.'' Holding out a coil while secretly supplying its measured basis is not unknown-source generalization. For unseen coils, use CAD-based physical templates or new independent arrays and disclose the information provided.

\caution{A known activation vector can label which source was driven; it does not necessarily give an exact current amplitude or a unique multilayer current map. Model-assisted labels, nominal heights, and independently measured targets should appear as separate evidence tiers.}
\deliver{Produce two manifests, a source/session split file, a dictionary reproduction report, and a list of missing metadata or acquisitions needed for the primary near-plane claim.}
